% Options for packages loaded elsewhere
\PassOptionsToPackage{unicode}{hyperref}
\PassOptionsToPackage{hyphens}{url}
\documentclass[
  10pt,
  a4paper,
]{article}
\usepackage{xcolor}
\usepackage[margin=22mm]{geometry}
\usepackage{amsmath,amssymb}
\setcounter{secnumdepth}{-\maxdimen} % remove section numbering
\usepackage{iftex}
\ifPDFTeX
  \usepackage[T1]{fontenc}
  \usepackage[utf8]{inputenc}
  \usepackage{textcomp} % provide euro and other symbols
\else % if luatex or xetex
  \usepackage{unicode-math} % this also loads fontspec
  \defaultfontfeatures{Scale=MatchLowercase}
  \defaultfontfeatures[\rmfamily]{Ligatures=TeX,Scale=1}
\fi
\usepackage{lmodern}
\ifPDFTeX\else
  % xetex/luatex font selection
\fi
% Use upquote if available, for straight quotes in verbatim environments
\IfFileExists{upquote.sty}{\usepackage{upquote}}{}
\IfFileExists{microtype.sty}{% use microtype if available
  \usepackage[]{microtype}
  \UseMicrotypeSet[protrusion]{basicmath} % disable protrusion for tt fonts
}{}
\makeatletter
\@ifundefined{KOMAClassName}{% if non-KOMA class
  \IfFileExists{parskip.sty}{%
    \usepackage{parskip}
  }{% else
    \setlength{\parindent}{0pt}
    \setlength{\parskip}{6pt plus 2pt minus 1pt}}
}{% if KOMA class
  \KOMAoptions{parskip=half}}
\makeatother
\setlength{\emergencystretch}{3em} % prevent overfull lines
\providecommand{\tightlist}{%
  \setlength{\itemsep}{0pt}\setlength{\parskip}{0pt}}
\usepackage{mathrsfs}
\newcommand{\Log}{\operatorname{Log}}
\usepackage{mathtools}
\usepackage{xurl}
\emergencystretch=3em
\setcounter{tocdepth}{1}
\newcommand{\Beta}{\mathrm{B}}
\DeclareUnicodeCharacter{00B2}{\ensuremath{^2}}
\DeclareUnicodeCharacter{00B1}{\ensuremath{\pm}}
\DeclareUnicodeCharacter{00D7}{\ensuremath{\times}}
\DeclareUnicodeCharacter{2192}{\ensuremath{\to}}
\DeclareUnicodeCharacter{21D2}{\ensuremath{\Rightarrow}}
\DeclareUnicodeCharacter{21A6}{\ensuremath{\mapsto}}
\DeclareUnicodeCharacter{2208}{\ensuremath{\in}}
\DeclareUnicodeCharacter{2209}{\ensuremath{\notin}}
\DeclareUnicodeCharacter{2212}{\ensuremath{-}}
\DeclareUnicodeCharacter{2227}{\ensuremath{\wedge}}
\DeclareUnicodeCharacter{2228}{\ensuremath{\vee}}
\DeclareUnicodeCharacter{2260}{\ensuremath{\ne}}
\DeclareUnicodeCharacter{2264}{\ensuremath{\le}}
\DeclareUnicodeCharacter{2265}{\ensuremath{\ge}}
\DeclareUnicodeCharacter{2297}{\ensuremath{\otimes}}
\DeclareUnicodeCharacter{00B3}{\ensuremath{^3}}
\DeclareUnicodeCharacter{2074}{\ensuremath{^4}}
\DeclareUnicodeCharacter{2076}{\ensuremath{^6}}
\DeclareUnicodeCharacter{03BB}{\ensuremath{\lambda}}
\DeclareUnicodeCharacter{03BC}{\ensuremath{\mu}}
\DeclareUnicodeCharacter{03B5}{\ensuremath{\epsilon}}
\DeclareUnicodeCharacter{03C4}{\ensuremath{\tau}}
\DeclareUnicodeCharacter{03A6}{\ensuremath{\Phi}}
\DeclareUnicodeCharacter{03B1}{\ensuremath{\alpha}}
\DeclareUnicodeCharacter{03C3}{\ensuremath{\sigma}}
\DeclareUnicodeCharacter{03B4}{\ensuremath{\delta}}
\DeclareUnicodeCharacter{03C0}{\ensuremath{\pi}}
\DeclareUnicodeCharacter{039B}{\ensuremath{\Lambda}}
\DeclareUnicodeCharacter{2202}{\ensuremath{\partial}}
\DeclareUnicodeCharacter{0393}{\ensuremath{\Gamma}}
\DeclareUnicodeCharacter{2205}{\ensuremath{\varnothing}}
\DeclareUnicodeCharacter{221E}{\ensuremath{\infty}}
\DeclareUnicodeCharacter{2200}{\ensuremath{\forall}}
\DeclareUnicodeCharacter{00B9}{\ensuremath{^1}}
\DeclareUnicodeCharacter{00B0}{\ensuremath{{}^\circ}}
\DeclareUnicodeCharacter{211D}{\ensuremath{\mathbb{R}}}
\usepackage{bookmark}
\IfFileExists{xurl.sty}{\usepackage{xurl}}{} % add URL line breaks if available
\urlstyle{same}
\hypersetup{
  hidelinks,
  pdfcreator={LaTeX via pandoc}}

\author{}
\date{}

\begin{document}

{
\setcounter{tocdepth}{3}
\tableofcontents
}
\section{Faithful completion, local fibres and the uncompleted zeta heat
equation}\label{faithful-completion-local-fibres-and-the-uncompleted-zeta-heat-equation}

This derivation retains the original meromorphic zeta function, the
complete completion factor, its factorization, every local fibre and
every support label. The completed function is a coordinate in an
explicitly isomorphic invertible sheaf. Ordinary evaluation after
forgetting that sheaf is a different map, computed below. The heat
family is the original Rodgers--Tao family in its original coordinates.
This derivation neither assumes an absolute-base origin for that author
construction nor asserts Riemann-hypothesis positivity.

\subsection{1. Sources, conventions and the original
family}\label{sources-conventions-and-the-original-family}

The original author source is Brad Rodgers and Terence Tao, \emph{The de
Bruijn--Newman constant is non-negative}, arXiv:1801.05914v5, equations
hoz, sas, phidef, htdef and the heat-equation sentence following htdef;
\href{https://arxiv.org/src/1801.05914v5}{author TeX}. Its canonical
local ID is PUBUNIT-8AD81C89CA3B6B4D23E6BC27; the author file
\texttt{sources/\allowbreak{}rodgers\_\allowbreak{}tao\_\allowbreak{}1801.\allowbreak{}05914v5/\allowbreak{}fmp-\allowbreak{}template.\allowbreak{}tex} has SHA256
\texttt{25b015c0\allowbreak{}fb151f16\allowbreak{}13247d82\allowbreak{}cf53a656\allowbreak{}1a320d4d\allowbreak{}ab2bb4f8\allowbreak{}09482e7c\allowbreak{}477a5817}.
The introductory definitions and their surrounding statements were read,
not inferred from a search hit.

The programme receivers used here are
\href{HEAT_CAUCHY_ARITHMETIC_DERIVATION.md}{HA1--16, HA20},
\href{ACTUAL_HEAT_ZERO_DISTRIBUTION_VARIATION.md}{AG1--9, AG18--24},
\href{REAL_TIME_HEAT_TRACE_DERIVATION.md}{RT1--18}, and
\href{SUPPORTED_ZERO_PRIME_WEIL_DERIVATION.md}{SZW1--19, SZW24--38}. HA
supplies the unchanged completion logarithmic derivative and full
arithmetic terms; AG supplies the rational trace theorem and specified
contour variation; RT supplies the actual right-time compact-test
derivative. SZW supplies the supported-zero theta comparison and label
spaces. This note proves the additional completion/sheaf comparisons and
local calculations in full.

Write \[
\begin{aligned}
B(s)&=\pi^{-s/2}\Gamma(s/2),&
p(s)&=\tfrac12s(s-1),&
C(s)&=p(s)B(s),\\
\xi(s)&=C(s)\zeta(s),&
\Lambda(s)&=B(s)\zeta(s),&
\xi(s)&=p(s)\Lambda(s).
\end{aligned}\tag{UZ1}
\] The symbol \(p(s)\) in this factorization is a polynomial, not a
finite prime. The factor \(1/2\), the two endpoint factors and the whole
Gamma factor remain in every map.

The original heat family and its transported coordinate are \[
\begin{aligned}
\Phi(u)&=\sum_{n\ge1}(2\pi^2n^4e^{9u}-3\pi n^2e^{5u})e^{-\pi n^2e^{4u}},\\
H_t(z)&=\int_0^\infty e^{tu^2}\Phi(u)\cos(zu)\,du,\\
\xi_t(s)&=8H_t(-2i(s-\tfrac12)),\qquad
g_t(s)=2\xi_t(s)=16H_t(-2i(s-\tfrac12)),\\
\zeta_t(s)&=C(s)^{-1}\xi_t(s),\qquad
\Lambda_t(s)=B(s)\zeta_t(s)=\frac{2\xi_t(s)}{s(s-1)}.
\end{aligned}\tag{UZ2}
\] Thus \(\xi_0=\xi,\ \zeta_0=\zeta,\ g_0=2\xi\), and
\(\Lambda_t=g_t/(s(s-1))\), exactly as in HA20.

For \(u\ge0\), splitting the exponential in the \(n\)-sum gives \[
|\Phi(u)|\le K e^{9u-ce^{4u}}\quad(c>0).
\tag{UZ3}
\] Indeed one half of \(\pi n^2e^{4u}\) bounds a positive multiple of
\(e^{4u}\), while the other half makes \(\sum n^4e^{-\pi n^2/2}\)
summable. On compact sets of complex \(s,t\), derivatives of the
integrand add powers of \(u\), still bounded by an integrable function
\(K'\exp(Tu^2+(2R+10)u-ce^{4u})\). Consequently \(\xi_t(s)\) is jointly
entire in \((t,s)\), and direct differentiation, retaining
\((-2i)^2=-4\), gives \[
\partial_t\xi_t=\tfrac14\partial_s^2\xi_t,\qquad
\xi_t(1-s)=\xi_t(s).
\tag{UZ4}
\] For no complex \(t\) is \(\xi_t\) the zero function. To see this,
extend \(e^{tu^2}\Phi(u)\) evenly from the positive half-line. Its
extension is in \(L^1(\mathbb R)\) by (UZ3), is nonzero, and its Fourier
transform is \(2H_t\) on the real axis. If that transform vanished
identically, convolution with each Gaussian approximate identity would
vanish by Fourier inversion for the Gaussian and Fubini's theorem.
Approximation in \(L^1\) would make the original function zero, a
contradiction. Thus all divisors used below are defined, including at
complex times.

For real \(t,s\), every summand of \(\Phi(u)\) is positive because
\(2\pi n^2e^{4u}-3>0\), and the cosine in (UZ2) becomes
\(\cosh(2(s-\tfrac12)u)>0\). Hence \[
\xi_t(s)>0\quad(t,s\in\mathbb R).
\tag{UZ5}
\] This does not assert positivity of a Weil quadratic form.

\subsection{2. Gamma factors and their exact local
units}\label{gamma-factors-and-their-exact-local-units}

Here are the Gamma facts and their sufficient derivation. Euler's limit
and its reciprocal product are \[
\Gamma(z)=\lim_{N\to\infty}\frac{N!\,N^z}{z(z+1)\cdots(z+N)},\qquad
\frac1{\Gamma(z)}
=z e^{\gamma z}\prod_{n\ge1}(1+z/n)e^{-z/n}.
\tag{UZ6}
\] The first follows for \(\Re z>0\) from \[
\int_0^1x^{z-1}(1-x)^Ndx
=\frac{N!}{z(z+1)\cdots(z+N)}
\] by repeated integration by parts, substitution \(y=Nx\), and
dominated convergence using \((1-y/N)^N\le e^{-y}\). For the reciprocal
product use \(\sum_{n\le N}1/n-\log N\to\gamma\). The factors after
their linear terms have uniformly summable logarithms on every compact
set avoiding \(-1,-2,\ldots\). This proves local uniform convergence and
nonvanishing there. Thus the reciprocal is entire, its zeros are simple
and occur exactly at \(0,-1,-2,\ldots\), and \(\Gamma\) has no zeros.
The integral also gives \(\Gamma(z+1)=z\Gamma(z)\) and \(\Gamma(1)=1\).

Taking logarithms near zero gives the convergent identity \[
\log\Gamma(1+z)
=-\gamma z+\sum_{k\ge2}\frac{(-1)^k\zeta(k)}k z^k,\qquad |z|<1.
\tag{UZ7}
\] To verify it, expand \(\log(1+z/n)\) in (UZ6), interchange the
absolutely convergent double series for \(|z|<1\), and retain the linear
term separately. In this formula \(\zeta(k)=\sum_{n\ge1}n^{-k}\), with
\(k\ge2\), needs no analytic continuation.

Let \(u=s-s_*\). Write \(C(s_*+u)=u^d c_{s_*}(u)\), where \(c_{s_*}\) is
a holomorphic unit. At zero there is no residual zero or pole: \[
d=0,\qquad
c_0(u)=(u-1)\pi^{-u/2}\Gamma(1+u/2),\qquad c_0(0)=-1,
\tag{UZ8}
\] and the full convergent unit expansion is specified by \[
\log\frac{c_0(u)}{-1}
=\log(1-u)-\frac{\gamma+\log\pi}{2}u
+\sum_{k\ge2}\frac{(-1)^k\zeta(k)}{k2^k}u^k,\qquad |u|<1.
\tag{UZ9}
\] The branch has value zero at \(u=0\). In particular
\(C'(0)=1+(\gamma+\log\pi)/2\).

At \(s=1\), \[
d=1,\qquad
c_1(u)=\tfrac12(1+u)\pi^{-(1+u)/2}\Gamma((1+u)/2),\qquad
c_1(0)=\tfrac12,
\tag{UZ10}
\] and \[
\begin{aligned}
\log\frac{c_1(u)}{1/2}
={}&\log(1+u)
-\left(\frac{\gamma+\log\pi}{2}+\log2\right)u\\
&+\sum_{k\ge2}\frac{(-1)^k(1-2^{-k})\zeta(k)}k u^k,
\qquad |u|<1.
\end{aligned}\tag{UZ11}
\] For the value \(\Gamma(1/2)=\sqrt\pi\), square the positive Gaussian
integral and change to polar coordinates. For the logarithmic series,
differentiating (UZ6) gives
\(\psi(z)=-\gamma+\sum_{n\ge0}(1/(n+1)-1/(n+z))\). At \(z=1/2\) the
partial sums reduce to \(2(H_{N+1}-H_{2N+2})\to-2\log2\), so
\(\psi(1/2)=-\gamma-2\log2\). For \(k\ge2\), termwise differentiation
gives the Taylor coefficient \((-1)^k\sum_{n\ge0}(n+1/2)^{-k}/k\); after
the substitution \(z=1/2+u/2\) it is \((-1)^k(1-2^{-k})\zeta(k)/k\).
This proves (UZ11), including its linear coefficient.

At \(s=-2m,\ m\ge1\), put \[
\kappa_m=(-1)^m\frac{2m(2m+1)\pi^m}{m!}.
\tag{UZ12}
\] The complete local expression, retaining its original factors, is \[
\begin{aligned}
d&=-1,\qquad C(-2m+u)=u^{-1}c_{-2m}(u),\\
c_{-2m}(u)
&=\kappa_m
\left(1-\frac{u}{2m}\right)
\left(1-\frac{u}{2m+1}\right)
\pi^{-u/2}\Gamma(1+u/2)
\prod_{j=1}^{m}\left(1-\frac{u}{2j}\right)^{-1}.
\end{aligned}\tag{UZ13}
\] Indeed Gamma recurrence gives \[
\Gamma(-m+u/2)
=\frac{2(-1)^m}{m!\,u}\Gamma(1+u/2)
\prod_{j=1}^m(1-u/(2j))^{-1}.
\] Multiplication by \(\tfrac12(-2m+u)(-2m-1+u)\pi^{m-u/2}\) proves
(UZ13). In particular \(c_{-2m}(0)=\kappa_m\), with its sign retained.

For \(H_m^{(k)}=\sum_{j=1}^m j^{-k}\) and \(H_m=H_m^{(1)}\), the full
convergent logarithmic series on \(|u|<1\) is \[
\begin{aligned}
\log\frac{c_{-2m}(u)}{\kappa_m}
={}&\left[-\frac1{2m}-\frac1{2m+1}
-\frac{\gamma+\log\pi}{2}+\frac{H_m}{2}\right]u\\
&+\sum_{k\ge2}\frac1k
\left[
\frac{(-1)^k\zeta(k)+H_m^{(k)}}{2^k}
-\frac1{(2m)^k}-\frac1{(2m+1)^k}
\right]u^k .
\end{aligned}\tag{UZ14}
\] It follows by expanding every factor in (UZ13); convergence is
absolute on compact subdiscs. Equations (UZ9), (UZ11), (UZ14), followed
by exponentiation, specify all unit coefficients, not only leading
values. At every other finite point \(s_*\), \(d=0\) and
\(c_{s_*}(u)=C(s_*+u)\) is already a holomorphic unit.

The full factor divisors are therefore \[
\operatorname{div}B=-[0]-\sum_{m\ge1}[-2m],\quad
\operatorname{div}p=[0]+[1],\quad
D:=\operatorname{div}C=[1]-\sum_{m\ge1}[-2m].
\tag{UZ15}
\] Each sum is locally finite. The \([0]\) terms cancel as divisors; the
factors \(B,p\), their residues and their full units remain part of
(UZ1). In particular ``zero divisor coefficient at 0'' does not say that
these factors or the supported-zero endpoint data were absent.

\subsection{3. The invertible sheaf and all local
jets}\label{the-invertible-sheaf-and-all-local-jets}

Let \(X=\mathbb C\), \(\mathcal O=\mathcal O_X\) be the holomorphic
sheaf and \(\mathcal M\) the meromorphic sheaf. Define the embedded
fractional ideal sheaf \[
\mathcal L=\mathcal O(D)
:=\{f\in\mathcal M:\operatorname{div}f+D\ge0\}
=C^{-1}\mathcal O\subset\mathcal M.
\tag{UZ16}
\] The inequality is local at every finite point. Multiplication is an
everywhere invertible sheaf morphism \[
M_C:\mathcal L\xrightarrow{\;\sim\;}\mathcal O,\quad f\mapsto Cf,
\qquad M_C^{-1}(F)=C^{-1}F .
\tag{UZ17}
\] Proof: both compositions are the identity in \(\mathcal M\); (UZ16)
is precisely the condition that \(Cf\) be holomorphic. At \(s_*\), if
\(C=u^dc(u)\), then \(\mathcal L_{s_*}=u^{-d}\mathcal O_{s_*}\), because
\(c\) is a unit. Thus the morphism is an isomorphism at every stalk, not
only off the completion divisor. Multiplication by \(C\) is also an
automorphism of the \(\mathcal O\)-module \(\mathcal M\), but retaining
\((\mathcal O,\mathcal L,\mathcal O\hookrightarrow\mathcal M,\mathcal L\hookrightarrow\mathcal M,C)\)
additionally retains the original orders of vanishing and poles.

Write \(f=u^{-d}y(u)\), \(Cf=x(u)=c(u)y(u)\), with
\(x=\sum x_j u^j,\ y=\sum y_j u^j,\ c=\sum c_j u^j\). For every
\(N\ge0\), (UZ17) induces the fibre-and-jet isomorphism \[
\begin{aligned}
\mathcal L_{s_*}/u^{N+1}\mathcal L_{s_*}
&\xrightarrow{\;\sim\;}
\mathcal O_{s_*}/u^{N+1}\mathcal O_{s_*},\\
x_j&=\sum_{i=0}^{j}c_i y_{j-i},\qquad
y_j=c_0^{-1}\left(x_j-\sum_{i=1}^{j}c_i y_{j-i}\right)
\quad(0\le j\le N).
\end{aligned}\tag{UZ18}
\] The diagonal coefficient is \(c_0\ne0\), so the displayed recurrence
proves bijectivity for every finite jet order and for the inverse-limit
formal completion. Analytic sections are related by the convergent
multiplication and division in (UZ17). This keeps all derivatives and
Laurent coefficients with their exact index shift.

At the three types of special point, the fibre maps are as follows: \[
\begin{array}{c|c|c|c}
s_* & \mathcal L_{s_*} & \mathcal L_{s_*}/u\mathcal L_{s_*}
& M_C\text{ on that fibre}\\ \hline
-2m & u\mathcal O_{s_*} & u\mathcal O_{s_*}/u^2\mathcal O_{s_*}
& [u y]\mapsto \kappa_m y(0)\\
1 & u^{-1}\mathcal O_1 & u^{-1}\mathcal O_1/\mathcal O_1
& [u^{-1}y]\mapsto \tfrac12 y(0)\\
0 & \mathcal O_0 & \mathcal O_0/u\mathcal O_0
& [y]\mapsto -y(0).
\end{array}\tag{UZ19}
\] At \(-2m\), the inclusion \(u\mathcal O\hookrightarrow\mathcal O\)
induces the zero map into \(\mathcal O/u\mathcal O\), since \(uy\)
evaluates to zero. Its source fibre \(u\mathcal O/u^2\mathcal O\) is one
dimensional and need not be zero. The completion map instead multiplies
that fibre by \(\kappa_m\), an isomorphism. At 1 there is no ordinary
holomorphic evaluation of \(u^{-1}y\); the correct fibre records its
residue \(y(0)\). Evaluating \(C(1)=0\) and treating \(u^{-1}y\) as a
scalar value would not be a defined map. These conclusions compute the
maps responsible for the apparent losses.

There is also the faithful intermediate factorization \[
\begin{aligned}
\mathcal L_\Lambda&=p^{-1}\mathcal O
=\mathcal O([0]+[1])\subset\mathcal M,\\
\mathcal L \xrightarrow{\,M_B\,}\mathcal L_\Lambda
&\xrightarrow{\,M_p\,}\mathcal O,\qquad
M_pM_B=M_C,
\end{aligned}\tag{UZ20}
\] with both arrows isomorphisms. In fact \(pf_\Lambda\) is holomorphic
iff \(f_\Lambda\in\mathcal L_\Lambda\), and \(p(Bf)=Cf\), which proves
the first arrow and its inverse. This preserves the Gamma pole at zero
before the \(s\)-factor acts.

For completeness, the classical values and leading coefficients follow
without assuming trivial-zero simplicity. On \(\Re s>1\), integrating
the counting function \(\lfloor x\rfloor\) gives \[
\zeta(s)=s\int_1^\infty \lfloor x\rfloor x^{-s-1}dx
=\frac{s}{s-1}-s\int_1^\infty\{x\}x^{-s-1}dx.
\tag{UZ21}
\] The last integral is holomorphic on \(\Re s>0\), by boundedness of
\(\{x\}\) and differentiation on compact subsets. Thus the residue of
\(\zeta\) at 1 is 1. Equations (UZ10) and \(\xi(1-s)=\xi(s)\) imply \[
\xi(1)=\xi(0)=\tfrac12,\quad
\zeta(0)=-\tfrac12,\quad
\operatorname{Res}_{s=0}\Lambda=-1,\quad
\operatorname{Res}_{s=1}\Lambda=1.
\tag{UZ22}
\] At \(-2m\), reflection gives \[
\begin{aligned}
\xi(-2m)&=\xi(2m+1)
=\frac{m(2m+1)(2m)!}{4^m m!\pi^m}\zeta(2m+1)>0,\\
\zeta'(-2m)&=\frac{\xi(-2m)}{\kappa_m}
=(-1)^m\frac{(2m)!}{2(2\pi)^{2m}}\zeta(2m+1).
\end{aligned}\tag{UZ23}
\] To obtain the first expression use
\(\Gamma(m+1/2)=(2m)!\sqrt\pi/(4^m m!)\), proved from Gamma recurrence
and the Gaussian integral. The second follows by the first-fibre map in
(UZ19). It is nonzero, so these are simple trivial zeros. No zero or
coefficient has been identified with the unsupported element.

For every real time, (UZ5) and the same local maps give the exact
persistent geometry \[
\begin{aligned}
\operatorname{Res}_{s=1}\zeta_t&=2\xi_t(1)>0,&
\zeta_t(0)&=-\xi_t(0)<0,\\
\zeta_t'(-2m)&=\xi_t(-2m)/\kappa_m,&
\operatorname{sgn}\zeta_t'(-2m)&=(-1)^m,\\
\operatorname{Res}_{s=0}\Lambda_t&=-2\xi_t(0),&
\operatorname{Res}_{s=1}\Lambda_t&=2\xi_t(1).
\end{aligned}\tag{UZ24}
\] In particular the pole and trivial-zero orders remain exactly simple
for real \(t\). For complex \(t\) the sheaf maps and divisor identity
remain valid, but a zero of \(\xi_t\) at one of these points can change
the total order; no fixed-order assertion is made for complex time.

\subsection{4. The complete transported heat
operator}\label{the-complete-transported-heat-operator}

Put \(q=C'/C\). As a meromorphic function its exact expression is \[
\begin{aligned}
q(s)&=\frac1s+\frac1{s-1}-\frac{\log\pi}{2}
+\frac12\psi(s/2),\\
q'(s)&=-\frac1{s^2}-\frac1{(s-1)^2}
+\frac14\psi^{(1)}(s/2),\\
\frac{C''}{C}(s)
&=\left(\frac1s+\frac1{s-1}-\frac{\log\pi}{2}
+\frac12\psi(s/2)\right)^2
-\frac1{s^2}-\frac1{(s-1)^2}
+\frac14\psi^{(1)}(s/2).
\end{aligned}\tag{UZ25}
\] Here \(\psi=\Gamma'/\Gamma\). These formulas are valid
meromorphically, with the removable singularity at 0 assigned the unit
value from (UZ9); for example \(q(0)=-1-(\gamma+\log\pi)/2\). At 1 and
\(-2m\) the residues of \(q\) are respectively \(+1\) and \(-1\),
precisely the divisor coefficients in (UZ15).

Since \(C\) does not depend on \(t\), (UZ4) proves \[
\boxed{\quad
\partial_t\zeta_t
=\tfrac14\left[
\partial_s^2\zeta_t
+2q(s)\partial_s\zeta_t
+\bigl(q'(s)+q(s)^2\bigr)\zeta_t
\right]
=\tfrac14 C^{-1}\partial_s^2(C\zeta_t).
\quad}\tag{UZ26}
\] It follows by expanding \(\partial_s^2(C\zeta_t)\), without deleting
any factor of \(C\). In particular the uncompleted function does not
satisfy the coefficient-free heat equation in this coordinate. For real
time, \(\zeta_t''\) has a pole of order three at 1 with nonzero leading
coefficient, by (UZ24), whereas \(\partial_t\zeta_t\) has at most a
simple pole there because \(C^{-1}\) is fixed and \(\partial_t\xi_t\) is
entire.

The operator is defined across all special points as the connection \[
\nabla_C:\mathcal L\to\mathcal L,\qquad
\nabla_C f=C^{-1}\partial_s(Cf),\qquad
\partial_t\zeta_t=\tfrac14\nabla_C^2\zeta_t .
\tag{UZ27}
\] It maps \(\mathcal L\) to itself because derivatives of a holomorphic
function are holomorphic. It satisfies \(\nabla_C(af)=a'f+a\nabla_C f\)
for \(a\in\mathcal O\). Thus its meromorphic coefficient presentation
does not make its action on the retained lattice undefined.

More explicitly, at every \(s_*\) write \(\zeta_t=u^{-d}y_t(u)\),
\(C=u^dc(u)\). The full local heat equation is \[
\partial_t y_t
=\tfrac14\left[y_t''+2\frac{c'}c y_t'
+\frac{c''}c y_t\right].
\tag{UZ28}
\] Indeed \(\xi_t=cy_t\), and the time-independent factor \(u^{-d}\)
cancels only as a change of local frame in (UZ26). Equivalently,
differentiating \(u^{-d}y\) and using \(q=d/u+c'/c\) cancels each
displayed \(1/u\) term in the connection. The complete unit \(c\), given
in (UZ8), (UZ10), (UZ13), remains. Its coefficients are holomorphic on
the local disc.

All jets evolve through these same maps. For
\(x_t=\xi_t=\sum x_j(t)u^j\), \[
4\dot x_j=(j+2)(j+1)x_{j+2}.
\tag{UZ29}
\] The dot denotes a complex analytic local time derivative, justified
by joint entire dependence in (UZ3). If \(a(u)=c'/c=\sum a_i u^i\),
\(b(u)=c''/c=\sum b_i u^i\), and \(y_t=\sum y_j(t)u^j\), (UZ28) gives \[
4\dot y_j=(j+2)(j+1)y_{j+2}
+2\sum_{i=0}^j a_i(j-i+1)y_{j-i+1}
+\sum_{i=0}^j b_i y_{j-i}.
\tag{UZ30}
\] The recurrence (UZ18) intertwines (UZ29) and (UZ30). Finite jets are
not falsely claimed to form an autonomous heat system: their derivatives
use the two additional jet orders displayed in these equations.

On \(\Re s>1\), at \(t=0\), the logarithmic derivative of \(\xi\) is \[
L_0=q+\zeta'/\zeta
=q-\sum_{n\ge2}\Lambda(n)n^{-s}.
\tag{UZ31}
\] This is exactly HA10, with its original endpoint and Gamma terms.
Differentiating the logarithmic derivative in time gives \[
\partial_t(\zeta_t'/\zeta_t)
=\partial_t(\xi_t'/\xi_t)
=\tfrac14\partial_s(\xi_t''/\xi_t).
\tag{UZ32}
\] The first equality holds because \(q\) is independent of \(t\). Hence
the full causal distribution in CH and its archimedean--prime mixed
terms are unchanged by the faithful uncompleted coordinate. In
particular replacing \(q\) by zero would change the actual heat
variation and would discard the mixed terms, rather than implement
(UZ17).

Reflection also transports exactly. Let \(R(F)(s)=F(1-s)\). Then \[
J_C:=M_C^{-1}RM_C,\qquad
(J_C f)(s)=\frac{C(1-s)}{C(s)}f(1-s),\qquad
J_C^2=\mathrm{id}_{\mathcal L}.
\tag{UZ33}
\] These are operators on global sections; on local sections the
reflected operator sends sections over \(U\) to sections over \(1-U\),
so its sheaf map covers the reflection of the base and is not
\(\mathcal O\)-linear over the identity map. Proof: \(Cf\) is
holomorphic, so its reflection is holomorphic on the reflected domain,
proving the stated codomain; composition cancels the two meromorphic
ratios and gives the identity. Also \[
\nabla_C J_C=-J_C\nabla_C,\qquad
\nabla_C^2J_C=J_C\nabla_C^2,\qquad
J_C\zeta_t=\zeta_t.
\tag{UZ34}
\] The first relation is the chain rule after conjugating by \(C\); the
others follow by a second application and (UZ4). No unweighted
reflection of the uncompleted meromorphic coordinate is substituted for
this operator.

\subsection{5. Divisors, finite contours and global
traces}\label{divisors-finite-contours-and-global-traces}

The exact meromorphic divisor relation for every time is \[
\operatorname{div}\zeta_t=\operatorname{div}\xi_t-D
=\operatorname{div}\xi_t+\sum_{m\ge1}[-2m]-[1].
\tag{UZ35}
\] It is an identity of locally finite signed divisors. It follows at
each stalk from
\(\operatorname{ord}(C\zeta_t)=\operatorname{ord}C+\operatorname{ord}\zeta_t\).
If \(\xi_t\) vanishes at a special point, the orders add there; the
formula still holds.

Let \(\Omega\) be bounded with a piecewise smooth boundary avoiding the
zeros and poles of the functions involved, and let \(A\) be holomorphic
near \(\overline\Omega\). The exact finite-contour map is \[
\begin{aligned}
\frac1{2\pi i}\int_{\partial\Omega}
 A(s)\frac{\zeta_t'(s)}{\zeta_t(s)}\,ds
&=\frac1{2\pi i}\int_{\partial\Omega}
 A(s)\frac{\xi_t'(s)}{\xi_t(s)}\,ds
-\frac1{2\pi i}\int_{\partial\Omega}A(s)q(s)\,ds\\
&=\sum_{\rho\in\Omega}\operatorname{ord}_\rho(\xi_t)A(\rho)
+\sum_{-2m\in\Omega}A(-2m)-1_{\{1\in\Omega\}}A(1).
\end{aligned}\tag{UZ36}
\] Each sum is finite and includes multiplicity. This is the residue
theorem applied to \(\zeta_t'/\zeta_t=\xi_t'/\xi_t-q\). It applies
without a summability assumption on an infinite trivial-zero sequence.
If a test has a pole inside \(\Omega\), its additional residues must
also be retained; (UZ36) is stated with the holomorphic-test hypothesis
to avoid suppressing them.

For a fixed such contour nonvanishing persists at sufficiently small
complex time, so differentiating (UZ36) gives equality of its two heat
variations, since the \(q\) integral is fixed. Thus the equality holds
before an infinite contour limit, including at every trivial-zero and
endpoint local geometry.

Now let \(A\) be rational, \(A(s)=O(s^{-2})\) at infinity, with poles
avoiding \(1,-2,-4,\ldots\) and the zeros of \(\xi_{t_*}\) at a real
time \(t_*\). There is a complex time neighbourhood of \(t_*\) in which
\[
\boxed{\quad
Z_{\zeta_t}^{\mathrm{div}}(A)
=Z_{\xi_t}(A)+\sum_{m\ge1}A(-2m)-A(1).
\quad}\tag{UZ37}
\] The left side is the sum of the signed divisor orders against \(A\);
the right side retains all its three parts. Both sums converge
absolutely. To prove it, (UZ3) yields the uniform growth
\(\log\max_{|s|\le R}|\xi_t(s)|\le K(R+1)\log(R+2)\): absorb \(Tu^2\)
into half of \(ce^{4u}\), then substitute \(x=e^{4u}\) in the remaining
integral and bound the resulting Gamma integral. Since
\(\xi_{t_*}(0)>0\), it remains nonzero in a small complex time disc.
Jensen's formula then gives \(N_t(R)=O(R\log(R+2))\), uniformly there.
Dyadic summation implies
\(\sum_{|\rho|>R}m_\rho|\rho|^{-2}=O(\log(R+2)/R)\). Also
\(A(-2m)=O(m^{-2})\). These bounds prove absolute convergence, and
summing (UZ35) proves (UZ37).

For clarity, the derivative in (UZ37) is a genuine derivative of the
infinite rational trace. The genus-one logarithmic derivative is \[
\frac{\xi_t'}{\xi_t}(s)
=b_t+\sum_\rho m_\rho\left(\frac1{s-\rho}+\frac1\rho\right).
\] It converges normally on compact sets avoiding the zeros by the
preceding square-summability estimate. Summing residues at the finite
pole set \(\mathcal P(A)\) gives \[
Z_{\xi_t}(A)
=-\sum_{a\in\mathcal P(A)}
\operatorname{Res}_{s=a}
\left(A(s)\frac{\xi_t'(s)}{\xi_t(s)}\right).
\tag{UZ38}
\] Indeed \(\sum_a\operatorname{Res}_a A=0\), so \(b_t\) and every
constant \(1/\rho\) correction contribute zero, and the rational
function \(A(s)/(s-\rho)\) has pole-set residues summing to
\(-A(\rho)\). Normal convergence permits termwise residues. The right
side is holomorphic in time while its finite pole values avoid zeros.
Hence \[
\partial_t Z_{\zeta_t}^{\mathrm{div}}(A)
=\partial_t Z_{\xi_t}(A)
=-\frac14\sum_{a\in\mathcal P(A)}
\operatorname{Res}_{s=a}
\left(A(s)\partial_s\frac{\xi_t''(s)}{\xi_t(s)}\right).
\tag{UZ39}
\] This proves equality of the actual derivatives, not an unsupported
interchange with individually moving zero sums.

For entire compact-test Mellin transforms this rational summability
conclusion must not be transplanted. With the original convention \[
A(s)=M_h(s)=\int_{\mathbb R}h(v)e^{-(s-1/2)v}dv,
\tag{UZ40}
\] choose a nonzero nonnegative \(h\in C_c^\infty((a,b))\) with
\(0<a<b\). Then \[
M_h(-2m)\ge e^{(2m+1/2)a}\int h(v)dv,
\tag{UZ41}
\] so \(\sum_mM_h(-2m)\) diverges. Even endpoint filtering does not
repair this: for \(T=\partial_v^2-1/4\), two integrations by parts give
\[
M_{Th}(s)=s(s-1)M_h(s),
\] so its values at \(-2m\) have the additional positive factor
\(2m(2m+1)\) and still diverge. The values at 0 and 1 vanish, which is a
different exact statement.

There is nevertheless a faithful finite-contour and regularized map.
Retain at every \(\Omega\) the pair \[
\left(
Z_{\zeta_t,\Omega}^{\mathrm{div}}(A),\
D_\Omega(A):=\frac1{2\pi i}\int_{\partial\Omega}Aq\,ds
\right),\qquad
Z_{\zeta_t,\Omega}^{\mathrm{div}}(A)+D_\Omega(A)=Z_{\xi_t,\Omega}(A).
\tag{UZ42}
\] Retain also the actual divisors \(\operatorname{div}\zeta_t|_\Omega\)
and \(D|_\Omega\), rather than only their value under a single test.
Restriction to nested domains then records each added divisor point and
its order; this is the faithful locally finite divisor data. A single
weighted value can vanish at a nonempty divisor and is not an injective
encoding. For real \(t\ge0\), RT1--18 provides the critical strip and
complete compact-test trace. Exhausting the plane by bounded domains
with boundaries avoiding the divisors then gives the specified relative
trace \[
Z_{\zeta_t;C}^{\mathrm{rel}}(M_h)
:=\lim_{\Omega\uparrow\mathbb C}
\left(Z_{\zeta_t,\Omega}^{\mathrm{div}}(M_h)+D_\Omega(M_h)\right)
=Z_{\xi_t}(M_h).
\tag{UZ43}
\] The limit exists because the terms inside parentheses equal the
finite completed trace exactly and the completed zero sum converges
absolutely: compact smooth Mellin tests decrease faster than every power
on the closed critical strip, and the zero count above applies. The two
separate terms in parentheses need not converge. The definition is
therefore a relative trace with its fixed divisor correction explicitly
retained, not a convergent unsigned trivial-zero sum.

At time zero, on the separated-height contours of AG22, or any
enlargement adding the fixed real divisor points and excluding zeros on
its boundary, the \(D_\Omega\) contribution has derivative zero. The
limit of finite-contour derivatives is consequently the original AG20
variation. RT16 further proves that on the real \(t\ge0\) family this
variation is the actual right derivative of (UZ43) for compact smooth
\(h\). These are the specified contour map and the independently proved
actual-time upgrade; no unproved complex-time entire-test convergence is
inserted.

\subsection{6. The supported carrier and retained endpoint
representations}\label{the-supported-carrier-and-retained-endpoint-representations}

Let \(L\) be the original finite bounded distributive support lattice.
Use exactly the SZW10 carrier, extended to meromorphic amplitudes: \[
G_L(\mathcal M)
=\{(f,1_L):f\in\mathcal M\}
\cup\{z_\lambda=(0,\lambda):\lambda\ne1_L\},\qquad
\tau=z_{0_L},\quad e=z_{1_L}.
\] No nonzero amplitude is assigned a nontop label. The original
addition uses amplitude addition and support join; multiplication uses
amplitude multiplication and support meet. These operations preserve the
displayed set: if an amplitude sum or product is nonzero, each factor
that could force a nontop support has been excluded by that nonzero
amplitude. More explicitly, a nonzero sum has at least one top-labelled
summand, and a nonzero product has both factors top-labelled. The
element \(\tau=z_{0_L}\) is the additive identity and is
multiplicatively absorbing. These operations respect restriction of
sections. In particular \(e\ne\tau\) for a nontrivial support lattice.
The multiplication map is \[
\widehat M_C(z_\lambda)=z_\lambda,\qquad
\widehat M_C(f,1_L)=(Cf,1_L),\qquad
\widehat M_C^{-1}(F,1_L)=(C^{-1}F,1_L).
\tag{UZ44}
\] It is an additive, label-preserving bijection between the carrier
over \(\mathcal L\) and the carrier over \(\mathcal O\); the two
prescriptions agree at \(e=(0,1_L)\). It is multiplication by
\((C,1_L)\) in the ambient meromorphic semiring, and its inverse is
multiplication by \((C^{-1},1_L)\). It fixes every \(z_\lambda\),
including the distinct elements \(e\) and \(\tau\).

For precision this is a module/scalar-multiplication isomorphism, not a
claim that multiplication by \(C\) is a ring homomorphism. Moreover
\(\mathcal O(D)\) is not generally closed under multiplication in its
embedding: its natural product is
\(\mathcal O(D)\otimes\mathcal O(D)\to\mathcal O(2D)\). The
corresponding exact multiplicative square uses \(M_{C^2}\) on
\(\mathcal O(2D)\), since \((Cf)(Cg)=C^2fg\). Thus product domains and
all labels are retained.

At a special point, first apply the stalk/fibre morphism (UZ18) on the
top amplitude, leaving each lower zero label unchanged. This gives, for
example, \[
([u y],1_L)\longmapsto(\kappa_m y(0),1_L)
\quad\text{at }-2m,\qquad
([u^{-1}y],1_L)\longmapsto(\tfrac12y(0),1_L)
\quad\text{at }1,\qquad z_\lambda\longmapsto z_\lambda.
\tag{UZ45}
\] If a top fibre coefficient is zero, its image is \(e=(0,1_L)\); it is
never \(\tau\). At \(-2m\), the ambient holomorphic evaluation instead
gives \(e\), regardless of the derivative coefficient. Equation (UZ19)
proves exactly why that different map loses the first nonzero
coefficient. The faithful sheaf map retains it, and the independent
lower labels remain fixed.

Finally SZW's full endpoint representation and every lower fixed
coordinate can be carried as direct, unchanged components: \[
M_C\oplus\mathrm{id}_{V_L}:
\mathcal L\oplus V_L\longrightarrow\mathcal O\oplus V_L,
\tag{UZ46}
\] where \(V_L\) is the original endpoint/lower-label vector space of
SZW13 and SZW33--38. This formula is read on the spaces of sections
together with the attached fixed representation; its inverse is
\(M_{C^{-1}}\oplus\mathrm{id}_{V_L}\). The arbitrary complex
\(c_\lambda\) are values of functions \textbf{on} the lower carrier
points, as in SZW11. They are not nonzero amplitudes \textbf{in} the
carrier at nontop labels. This distinction retains both the exact set in
SZW10 and all its independent function coordinates. Equations (UZ20),
(UZ22), (UZ24) additionally retain the meromorphic residues of
\(\Lambda_t\) at both endpoints. A fixed supported-zero trace, an
endpoint residue, a divisor order, a local section coefficient and an
unsupported element have the connecting maps above; none has been
replaced by another object merely because a scalar coefficient vanishes.

The classical author construction (UZ2) is presented over complex \(s\),
complex-valued functions and real heat time, with joint complex analytic
extension justified here. The programme's lift (UZ44) and its augmented
endpoint map (UZ46) are explicit additional constructions. Their
existence does not prove that the original author heat kernel was
defined over the absolute base. Completion is useful here because it
gives an entire heat coordinate and its reflection; uncompletion through
(UZ17), (UZ20), (UZ26), (UZ35) and (UZ42) preserves the full local and
supported data. The signs in (UZ23)--(UZ24) concern exact real local
coefficients, and do not establish positivity of the full reflected Weil
pairing.

\subsection{7. All reflected local orientations and the unchanged
nontrivial
divisor}\label{all-reflected-local-orientations-and-the-unchanged-nontrivial-divisor}

Write \(\mathcal R_C(s)=C(1-s)/C(s)\), retaining this quotient as the
definition. Substitution of (UZ1), with both polynomial factors retained
before the equality is taken, proves the meromorphic identity \[
\mathcal R_C(s)=
\frac{\frac12(1-s)(-s)\pi^{-(1-s)/2}\Gamma((1-s)/2)}
{\frac12s(s-1)\pi^{-s/2}\Gamma(s/2)}
=\pi^{s-1/2}\frac{\Gamma((1-s)/2)}{\Gamma(s/2)}.
\tag{UZ47}
\] The equality of the polynomial factors is an identity of meromorphic
functions, and does not remove the embedded lattices or their endpoint
fibres. By (UZ15), or directly by (UZ6), \[
\operatorname{div}\mathcal R_C
=\sum_{m\ge0}[-2m]-\sum_{m\ge0}[1+2m].
\tag{UZ48}
\] For \(m\ge1\), put \(r_m=(-1)^m(2m)!/[2(2\pi)^{2m}]\). The exact
leading orientations are \[
\begin{aligned}
\mathcal R_C(u)&=\tfrac12u+O(u^2),&
\mathcal R_C(1+u)&=-2u^{-1}+O(1),\\
\mathcal R_C(-2m+u)&=r_mu+O(u^2),&
\mathcal R_C(1+2m+u)&=-r_m^{-1}u^{-1}+O(1).
\end{aligned}\tag{UZ49}
\] At zero and one these follow from \(c_0(0)=-1,\ c_1(0)=1/2\), using
\(1-(1+u)=-u\). At \(-2m\), the leading quotient is
\(C(1+2m)/\kappa_m\), which equals \(r_m\) by the same Gamma recurrence
used in (UZ23). At the reflected point the local variable is \(-u\);
this accounts for the minus sign in the pole coefficient. Thus a sign is
not silently lost on reflection.

For an arbitrary point \(p_*\), set \(p_*^\vee=1-p_*\),
\(d_p=\operatorname{ord}_p C\), and \(C(p+u)=u^{d_p}c_p(u)\). Then \[
\mathcal R_C(p_*+u)
=(-1)^{d_{p_*^\vee}}u^{d_{p_*^\vee}-d_{p_*}}
\frac{c_{p_*^\vee}(-u)}{c_{p_*}(u)}.
\tag{UZ50}
\] This follows by factoring numerator and denominator at their
indicated points. It is valid also when \(d=-1\). Write
\(\zeta_t(p+u)=u^{-d_p}y_{p,t}(u)\). Substitution in
\(\zeta_t(s)=\mathcal R_C(s)\zeta_t(1-s)\) gives the exact regular-germ
and all-jet identity \[
y_{p_*,t}(u)
=\frac{c_{p_*^\vee}(-u)}{c_{p_*}(u)}\,y_{p_*^\vee,t}(-u).
\tag{UZ51}
\] The two factors \((-1)^{d_{p_*^\vee}}\) and \((-1)^{-d_{p_*^\vee}}\)
cancel in this frame map. Every \(j\)-th Taylor coefficient on the
reflected side still carries the orientation \((-1)^j\), and
multiplication/division by the displayed full unit supplies the
remaining exact triangular coefficients as in (UZ18). In particular at
\(0,1\), \(y_{0,t}(0)=-y_{1,t}(0)/2\); it is the relation between the
original value at zero and the original pole residue.

Complex conjugation can be retained too. Since the defining factors of
\(C\) are real under conjugation, define \[
J_C^\# f(s)=C(s)^{-1}
\overline{C(1-\bar s)f(1-\bar s)}
=\frac{C(1-s)}{C(s)}\,\overline{f(1-\bar s)}.
\tag{UZ52}
\] This is an antilinear involution on global sections of
\(\mathcal L\), covering \(s\mapsto1-\bar s\) locally; its square is the
identity by the same quotient cancellation. For real \(t\),
\(J_C^\#\zeta_t=\zeta_t\) follows from the integral (UZ2). The
comparison conjugates the reflected section involution exactly. It does
not replace the original test-function involution or assign a sign to
its trace pairing.

At a zero \(\rho\) of \(\xi_t\) away from the fixed completion divisor,
\(C\) is a holomorphic unit. Hence the local ideals generated by
\(\zeta_t\) and \(\xi_t\) coincide: \[
(\zeta_t)_\rho=(\xi_t)_\rho,\qquad
\mathcal O_\rho/(\zeta_t)_\rho
=\mathcal O_\rho/(\xi_t)_\rho.
\tag{UZ53}
\] This follows because each generator is the other multiplied by a
holomorphic unit. Thus the full nontrivial local divisor algebra,
multiplicity and nilpotent jet directions are preserved exactly.
Together with the distinct fractional fibres (UZ19) at the fixed special
points, this proves the complete local comparison rather than only a
correspondence of zero sets off those points.

\subsection{Original-zeta Gaussian and affine receiving
maps}\label{original-zeta-gaussian-and-affine-receiving-maps}

The complete extension OZG1--55 proves that Gaussian convolution of the
test at every positive time makes the original trivial-zero scalar sum
diverge for every real translation. Its admissible replacement retains
each cutoff trace and the Gamma boundary with the identical finite sum
U\_N; OZG20--27 gives both maps and the inverse. The full coefficient
tower, its raw non-Hermitian correction and exact compensated negative
index are OZG28--39. No compact-support prime window is assumed after
smoothing; all prime powers and Gamma terms, with explicit error bounds,
are OZG40--52. This extension acts on the tests while fixing the
original zeta zeros.

The different affine change of the actual original heat family is
OZR1--36. Its multiplier is exp(chi) C(phi)/C, with the full inverse,
exceptional units and local jets. The original generator contains every
q derivative term. The signed divisor comparison is OZR19--24; its
Cauchy logarithmic kernel has an additional growth contribution 2a.
OZR33 proves the exact full negative index after that contribution is
retained. The full support carrier and every lower coordinate are
OZG53--55 and OZR34--36. These are the receiving maps for these
specified extensions; the original preceding test, time and domains
remain those stated in its proof.

\subsection{Original endpoint sectors and the actual prime-boundary
return}\label{original-endpoint-sectors-and-the-actual-prime-boundary-return}

The complete new receiving proofs are
\href{SECTORIAL_ENDPOINT_ZERO_TRACE_RETURN.md}{SER1--27},
\href{SECTORIAL_FILTER_TRACE_DERIVATION.md}{SFT1--48}, and
\href{GAUSSIAN_ARITHMETIC_BOUNDARY_RETURN.md}{GAB}. They distinguish the
actual original parameters: SER/SFT uses the original spectral heat time
t; GAB uses the integer Gaussian coefficient d=1/(4 pi T), whose zero
limit is original T tending to infinity. Both full incoming programme
TeX sources are retained in supporting\_proofs with their human
references.

For the original unfiltered theta continuation, SER8 retains the full
Gamma multiplier and both exponential endpoint sectors. Its common-part
return has all right time derivatives but zero time-Taylor radius, with
exact remainder SER14. The full differential inverse is SFT12--22,
including one initial datum in the actual even domain. The zero divisors
change by the explicitly evaluated quotient SFT23--30, not by an
assumption based on kernel dimension. SER25 computes every
finite-contour first trace change, with all multiplicities. SFT39--48
proves a global cubic-decay rational trace, every Gamma contribution,
both endpoints and test/divisor collisions; it also evaluates a
quadratic-growth counterexample to dropping the infinity term. Thus the
original preceding filtered heat, compact-test and Cauchy calculations
retain their stated domains and do not silently become raw-sectorial
trace claims.

For the actual two-moment prime receiver, GAB10--16 evaluates the full
Mellin--Barnes formula, the singular Gamma sector and each positive-odd
logarithmic collision. GAB18--22 retains the inverse pair and every
additional pole. GAB32--44 puts back the full Mellin multiplier, marked
integer-zero source, actual supported-zero action and fixed prime label.
The first moment vanishes throughout the critical strip, while its
retained residual has exact order d\^{}(3/2) precisely on the original
nontrivial zero divisor. The coefficient there is nonzero, proved using
original zeta at rho−2. This is an actual all-zero detector, not a sign
assigned to the full Weil pairing. No new seminorm or scalar base is
substituted.

\end{document}
